\section{Ontología HMT--MD de las partículas}
\Needspace{7\baselineskip}
\subsection{Estado genealógico material: extensión, ruta, registro, firma y representación}
Un estado HMT completo no es un punto desnudo. Es la tupla

\[
 \mathfrak x=
 (\Gamma,\mathcal L_\Gamma,\sigma_\Gamma,\eta_\Gamma,
 \mathcal R_\Gamma,\mathcal B_\Gamma,\mathcal C_\Gamma),
\]
donde \(\Gamma\) es una ruta superviviente, \(\mathcal L_\Gamma\) su registro genealógico, \(\sigma_\Gamma\) su firma entera, \(\eta_\Gamma\) su refinamiento armónico, \(\mathcal R_\Gamma\) su representación, \(\mathcal B_\Gamma\) su memoria de retorno y \(\mathcal C_\Gamma\) su dato de composición o frontera. Dos estados con igual palabra visible pueden ser distintos si difieren en prefijo, bloque, memoria, acarreo, hoja o sombra de Witt.

\Needspace{7\baselineskip}
\subsection{Partícula como sección local persistente del estado genealógico dimensional}
Una partícula es una clase de equivalencia de secciones persistentes que satisface cuatro condiciones:

\begin{enumerate}
\item prolongación compatible a través de los cilindros del TPK;
\item retorno de fase con memoria de estado;
\item representación física irreducible o bloque mínimo invariante;
\item estabilidad espectral bajo el acoplamiento que la observa.
\end{enumerate}
Una sección persistente es una familia compatible

\[
 s=(s_K)_{K\ge K_0},\qquad
 p_{L,K}(s_L)=s_K\quad(L\ge K\ge K_0).
\]
Dos secciones \(s\) y \(s'\) representan la misma partícula cuando existe un nivel \(K_1\) y una familia compatible de simetrías internas \(g_K\in G_K\) tal que

\[
 s'_K=g_Ks_K\quad(K\ge K_1),
 \qquad
 p_{L,K}g_L=g_Kp_{L,K},
\]
y los caracteres observables del registro genealógico coinciden. Esta relación identifica descripciones que difieren por calibre interno, pero no identifica rutas con memoria, acarreo, hoja, bloque o holonomía distintos. La partícula es la clase

\[
 [s]_{\rm pers}
 =\{s':s'\sim s\},
\]
no una palabra visible aislada ni una celda elegida después de conocer su masa.

Si \(\widetilde\Gamma_K\) es la extensión al nivel \(K\), la persistencia se expresa por

\[
 p_{K+9,K}(\widetilde\Gamma_{K+9})
 =\widetilde\Gamma_K,\qquad
 g(K+9)=g(K),
\]
sin exigir

\[
 \mathcal B_{K+9}=\mathcal B_K.
\]
La partícula no es la fase que retorna, sino la identidad genealógica conservada a través del cambio de estado. Su masa mide la rigidez espectral de esa persistencia cuando se acopla a una realización material.

\Needspace{7\baselineskip}
\subsection{Antipartícula como imagen conjugada de una sección persistente bajo inversión de hoja}
La conjugación dodecafásica es

\[
 C_M(\mathbf t)=-\operatorname{rev}(\mathbf t).
\]
El par partícula--antipartícula está formado por dos orientaciones de una misma banda genealógica. Si \(\chi\) denota la orientación de hoja,

\[
 (\mathbf t,\chi)\sim(C_M\mathbf t,-\chi).
\]
Para una energía de banda

\[
 E_{\rm band}(\mathbf t,\chi)
 =E_+(\mathbf t)+\chi E_-(\mathbf t),
\]
con

\[
 E_\pm(\mathbf t)
 =\frac12\bigl(E_0(\mathbf t)\pm E_0(C_M\mathbf t)\bigr),
\]
se obtiene

\[
 E_{\rm band}(C_M\mathbf t,-\chi)
 =E_{\rm band}(\mathbf t,\chi).
\]
La igualdad de masas es, por tanto, consecuencia de la paridad del operador bajo la conjugación, mientras que carga y orientación cambian. La igualdad de anchuras requiere además que el operador de decaimiento sea conjugado por \(C_M\).

\Needspace{7\baselineskip}
\subsection{Masa como carácter espectral de persistencia de \(1\) ruta material}
La masa es la medida espectral, en la recta \(\mathcal L_M\), del operador que resulta de comprimir carácter, torres, lector bidireccional y reloj sobre el subespacio persistente seleccionado por el acoplamiento:

\[
 \widehat M_{P,a}
 =P_{P,a}\,
 \mathbf m_\star^{\rm ret}
 R_{\rm act}^{\widehat K_{P,a}/2}
 \widehat{\mathcal A}_{P,a}^{(0)}
 R_{\rm act}^{\widehat K_{P,a}/2}
 P_{P,a}.
\]
La salida general es

\[
 \mu_{P,a}^{\rm mass}(B)
 =\operatorname{Tr}\!\left(
 \rho_{P,a}E_{\widehat M_{P,a}}(B)\right).
\]
Existe una masa escalar \(m_P\) si y sólo si

\[
 P_{\operatorname{supp}\rho}\widehat M_{P,a}
 P_{\operatorname{supp}\rho}
 =m_P P_{\operatorname{supp}\rho},
\]
equivalentemente, si la varianza de \(\widehat M_{P,a}\) en el estado preparado es cero. Esta definición incluye el electrón de rango uno, los multipletes forzados por Schur y los polos aislados de resonancias, sin confundirlos.

\Needspace{7\baselineskip}
\subsection{Carga como carácter de calibre transportado por la sección material}
La carga es el carácter orientado de la capa impar del registro genealógico bajo conjugación. Sea \(\sigma_{\rm odd}\) la firma impar. Entonces

\[
 Q(C_M\Gamma)=-Q(\Gamma),\qquad
 Q(\Gamma)=\mathfrak q\!\left(\sigma_{\rm odd}(\Gamma)\right),
\]
donde \(\mathfrak q\) es el homomorfismo de la firma impar al retículo de carga de la realización. La firma másica par y la firma de carga impar son objetos distintos: una palabra puede tener suma lineal nula y, sin embargo, portar una firma másica no nula.

\Needspace{7\baselineskip}
\subsection{Espín como representación de la acción rotacional sobre la fibra material}
El espín es la clase de holonomía con la que una sección persistente responde a una rotación completa del marco de orientación. La raíz interna

\[
 J^2=-I
\]
separa los dos cierres:

\[
 U(2\pi)=
 \begin{cases}
 +I,&\text{sector tensorial},\\
 -I,&\text{sector espinorial},
 \end{cases}
 \qquad
 U(4\pi)=I.
\]
Un bosón entero vive en una representación que cierra a \(2\pi\); un fermión semientero necesita la doble vuelta \(4\pi\). El valor de espín no se deduce de un conteo cardinal aislado: es el peso de la representación del grupo de holonomía sobre el subespacio persistente. En particular:

\[
 s=0:\text{ sección escalar},\qquad
 s=\tfrac12:\text{ banda espinorial},
\]
\[
 s=1:\text{ sección vectorial transversal},\qquad
 s=2:\text{ componente tensorial de una realización métrica}.
\]
El último renglón define un tipo de representación; no obliga a que exista una partícula elemental de espín dos.

La quiralidad y la helicidad son datos posteriores y distintos. Una involución orientada \(\Gamma_{\rm ch}^2=I\) define

\[
 P_L=\frac{I-\Gamma_{\rm ch}}2,
 \qquad
 P_R=\frac{I+\Gamma_{\rm ch}}2,
\]
mientras que la helicidad es el signo espectral del generador de rotación sobre la dirección de propagación. La primera clasifica hojas de la representación; la segunda depende del estado cinemático. Esta separación permite que un neutrino sea fermiónico por su cierre \(4\pi\), aunque su acoplamiento débil seleccione una sola quiralidad.

\Needspace{7\baselineskip}
\subsection{Estadística cuántica inducida por la simetría de intercambio de secciones}
La estadística es la regla de completación del espacio de rutas idénticas:

\[
 \mathcal F_-(\mathcal H)
 =\bigoplus_{n\ge0}\wedge^n\mathcal H,\qquad
 \mathcal F_+(\mathcal H)
 =\bigoplus_{n\ge0}\operatorname{Sym}^n\mathcal H.
\]
En la completación exterior, las relaciones CAR

\[
 \{a_i,a_j^\dagger\}=\delta_{ij},\qquad
 \{a_i,a_j\}=0
\]
implican

\[
 n_i^2=n_i
\]
y, por tanto, exclusión de Pauli y distribución de Fermi--Dirac. En la completación simétrica, las relaciones CCR permiten ocupación múltiple y conducen a Bose--Einstein. La estadística no se infiere de la masa: procede de cómo se compone la ruta bajo intercambio.

\Needspace{7\baselineskip}
\subsection{Color como representación residual de \(SU(3)\) sobre las rutas quark}
El color es la órbita residual triádica de una sección quark dentro de la fibra de carga y sabor. Las treinta y seis posiciones coloreadas se descomponen como

\[
 12_R+12_G+12_B.
\]
El objeto físico no es uno de esos colores, sino el triplete y su acción:

\[
 \mathcal H_q^{\rm color}
 =\mathcal H_R\oplus\mathcal H_G\oplus\mathcal H_B.
\]
Si el operador de masa conmuta con la acción irreducible de color,

\[
 U_g\widehat M_qU_g^*=\widehat M_q,
\]
el lema de Schur fuerza una masa común en el triplete. Elegir una celda roja, verde o azul como representante destruiría la covariancia.

\Needspace{7\baselineskip}
\subsection{Sabor como clase de ruta dentro de una generación material}
El sabor es una clase de acoplamiento inequivalente dentro de una misma representación de calibre, caracterizada por carga, generación, hoja, firma fina, memoria y proyector de interacción. Dos sabores pueden compartir una fibra gruesa y diferir en el entrelazador que los acopla al lector.

Formalmente, si \(F\) es la fibra común, los sabores \(f_i\) son representaciones \(\pi_i\) para las que

\[
 \operatorname{Hom}_{G_F}(V_{f_i},\ell^2(F))\ne0
\]
y cuyos proyectores \(P_{f_i}\) no son conjugados por una simetría de calibre que se mantenga física. El sabor no es la profundidad del atlas y no se reduce a un nombre convencional.

\Needspace{7\baselineskip}
\subsection{Generación material como estrato recurrente tipado por profundidad de retorno, firma fina y espectro del lector}
Una generación es un estrato recurrente de la misma gramática de carga y representación, distinguido por profundidad de retorno, firma fina y espectro del lector. La etiqueta de generación sólo es física cuando existe un proyector que entrelaza ese estrato con el canal de acoplamiento. Por ello, una profundidad \(G_j\) no debe identificarse automáticamente con electrón, muón, tau o neutrino estéril.

\Needspace{7\baselineskip}
\subsection{Campo de calibre como conexión sobre las fibras de carga}
Un campo de calibre es la realización de una redundancia de marco sobre las fibras. Sus bosones elementales pertenecen a la dirección nula del operador de masa mientras la simetría permanezca exacta. Una masa de calibre no se obtiene aproximando a cero un carácter positivo, sino mediante una ruptura, un acoplamiento longitudinal o una realización de polo declarada.

\Needspace{7\baselineskip}
\subsection{Partícula compuesta como sección persistente del producto de rutas enlazadas}
Una partícula compuesta es una sección persistente del operador de composición, no la suma de las masas ni de las firmas de sus constituyentes. Si \(\Gamma_1,\ldots,\Gamma_n\) son rutas constituyentes, primero se construye

\[
 \mathfrak C_{12}
 (\Gamma_1,\ldots,\Gamma_n;\mathcal T,\partial,\mathrm{carry})
 =\mathcal L_{\rm comp},
\]
donde \(\mathcal T\) es el árbol temporal y \(\partial\) la frontera de enlace. Sólo después actúa el lector

\[
 \mathcal L_{\rm comp}\longrightarrow\sigma_{\rm comp}
 \longrightarrow\chi_{\rm comp}
 \longrightarrow\widehat M_{\rm comp}.
\]
La energía de ligadura reside en el cociclo de composición, la memoria de frontera y la modificación del espectro; no se introduce como corrección externa.

\Needspace{7\baselineskip}
\subsection{Resonancia como sección transitoria y anchura como tasa de desintegración}
Una resonancia es un modo persistente no unitario del operador acoplado. Si

\[
 \lambda=re^{-i\theta},\qquad 0<r<1,
\]
entonces

\[
 E=\frac{\hbar_{\rm int}}{t_0}
 \bigl(\theta-i\log r\bigr)
 =M-\frac{i}{2}\Gamma.
\]
La masa \(M\) y la anchura \(\Gamma\) son las partes real e imaginaria de un polo, pero requieren lectores distintos. La anchura también puede obtenerse de un semigrupo comprimido:

\[
 T_X=Q_XUQ_X,\qquad
 \lambda_X=-\frac1{t_0}\log\rho(T_X),\qquad
 \Gamma_X=\hbar_{\rm int}\lambda_X.
\]
No se deduce una anchura únicamente de que una ruta posea horizonte finito.

\Needspace{7\baselineskip}
\subsection{Sección virtual como extensión intermedia no asintótica del atlas}
Una sección es virtual cuando su prolongación termina de manera no trivial:

\[
 \operatorname{Ext}_{L}(s)\ne\varnothing,\qquad
 \operatorname{Ext}_{L+1}(s)=\varnothing,\qquad
 \mathcal L(s)=L<\infty.
\]
Su contenido primario es el horizonte \(L\), la obstrucción terminal, el registro genealógico finito y la amplitud de propagación. No es una partícula de masa imaginaria ni un estado documental todavía no clasificado.

\Needspace{7\baselineskip}
